\section{Discussion}
\label{sec:discussion}

\subsection{Principal findings}
Under a validation design in which every data-dependent step was re-fitted inside each split, combining RFECV with an eight-learner stacking ensemble did not improve discrimination over regularised logistic regression on the PIMA cohort. Across 50 cross-validation folds, the stacks and logistic regression differed by less than 0.005 in ROC-AUC, and no difference approached significance after accounting for correlated folds and multiple comparisons. The same held on the hold-out set, where every DeLong comparison was non-significant. The higher hold-out sensitivity, F1 and MCC of RFECV + stacking (0.815, 0.710 and 0.532) are best read as the highest observed values on one 154-patient split, not as evidence of superiority. Its bootstrap confidence intervals overlap those of every comparator.

This null result is itself informative. It agrees with the systematic finding that flexible ML methods offer no average discrimination benefit over logistic regression in clinical prediction~\cite{christodoulou2019}. It also contrasts with the very high accuracies (0.90--0.98) sometimes reported for complex pipelines on the same 768 patients~\cite{zhou2023boruta,shams2025rfegru}. We cannot reproduce or audit those studies. However, a dataset with eight routine variables, a 35\% prevalence and a cross-validated ROC-AUC ceiling near 0.85 leaves little room for such gains. Validation design, especially preprocessing and feature selection outside the resampling loop~\cite{ambroise2002selection,kapoor2023leakage}, should be the first explanation considered for them.

\subsection{What RFECV and stacking contribute}
RFECV reduced the inputs by about a quarter with no loss of discrimination and identified a stable core of glycaemic, adiposity and family-history information, consistent with the SHAP analysis. The engineered interactions (\texttt{Glucose\_BMI}, \texttt{Glucose\_Age}, \texttt{BMI\_Age}, \texttt{HOMA\_IR\_surrogate}, \texttt{Pedigree\_Age}) were consistently preferred over thresholded flags and over raw blood pressure or skin-fold thickness. Two caveats apply. First, the selected size varied widely between folds (7--16 features) because the RFECV score curve is flat above about six features. A one-standard-error rule or stability selection~\cite{meinshausen2010stability} would yield smaller and more reproducible subsets. Second, the subset depended on the ranking estimator: coefficient-ranked RFECV for logistic regression chose different features than random-forest-ranked RFECV. The common protocol used here ensures a fair, identical selection procedure across learners, but it is not a learner-specific optimum, and results for individual learners should be read with that in mind.

Stacking added robustness rather than discrimination. It was never significantly worse than the best single learner, it did not require choosing a learner in advance, and at the default threshold it operated at higher sensitivity. Its meta-learner can also expose redundant base learners. Against these advantages stand higher computational cost (one RFECV + stacking fit took about 45~s, compared with milliseconds for logistic regression) and reduced transparency. On cohorts of this size, logistic regression remains a strong and simpler default.

\subsection{Calibration and thresholds matter more than the choice of learner}
The largest practical difference between configurations was not discrimination but calibration. Balanced class weights, which are widely used to raise sensitivity on imbalanced data, shifted predicted risks upwards by almost exactly the logit of the prevalence. The calibration slope stayed near one. The ranking of patients is therefore unaffected, but the predicted probabilities cannot be read as absolute risks without recalibration~\cite{vancalster2019calibration}. A simple intercept correction or Platt rescaling without class weights, fitted on out-of-fold predictions, would remove this offset. The well-calibrated SVM in our comparison shows the effect directly. Reporting only discrimination metrics, as most PIMA studies do, would have hidden this issue entirely.

Similarly, the default 0.5 threshold is arbitrary for screening. Thresholds chosen on development data for 80\%--90\% sensitivity transferred to the hold-out set within sampling error, at a cost of referring roughly half to 60\% of patients for confirmatory testing. Decision curves showed positive net benefit over the clinically plausible threshold range. However, all three models gave nearly identical net benefit, so the operating point matters more for clinical use than the choice between them.

\subsection{Missing data}
Almost half of the insulin values and 30\% of the skin-fold values were missing, so missingness itself might be expected to carry predictive information~\cite{sperrin2020missing}. In this cohort it did not. Missingness indicators slightly reduced discrimination, and iterative or $k$-NN imputation gave the same results as the simpler median imputation. The low SHAP contributions of insulin and skin-fold thickness agree with this. Whatever information these partially observed variables carry is largely redundant with glucose and BMI. Median imputation fitted within training folds was therefore adequate for this PIMA analysis under the evaluated alternatives; this may not hold in other cohorts, particularly where tests are ordered selectively.


\subsection{Limitations}
\begin{itemize}
  \item \textbf{One small, specific cohort.} PIMA comprises 768 women of one ethnic group from a single study. The results may not transfer to men, other populations or contemporary diagnostic criteria based on HbA1c. With 268 events and 16 candidate predictors, the cohort is also small for flexible model development~\cite{riley2020samplesize}.
  \item \textbf{No external validation.} Repeated cross-validation and a hold-out set estimate internal performance only. Validation on independent cohorts, such as population surveys with OGTT measurements, is needed before any clinical use.
  \item \textbf{One hold-out split.} With 54 positive hold-out patients, confidence intervals are wide. We therefore base inference on the 50-fold repeated cross-validation and treat hold-out results as confirmatory.
  \item \textbf{Fixed hyperparameters.} Learners were not tuned, to avoid an additional selection loop. Nested tuning could improve individual learners slightly, although the cross-validated ceiling suggests limited headroom.
  \item \textbf{Post-load measurements.} PIMA glucose and insulin are 2-h OGTT values, so \texttt{HOMA\_IR\_surrogate} is not a clinical HOMA-IR, and the model applies to settings where an OGTT has already been performed.
\end{itemize}

\subsection{Future work}
Future work should (i) validate the pipeline externally on larger, more diverse cohorts; (ii) incorporate recalibration and threshold selection into the nested pipeline; (iii) compare RFECV with stability selection and with learner-specific nested RFECV; and (iv) assess fairness and transportability across sex and ethnicity, as recommended by TRIPOD+AI~\cite{collins2024tripodai}.
